\chapter{Unpublished Results and Future Work}
\label{ch:future-work}

\lettrine{T}{he} preceding chapters presented published (or submitted) work.
This chapter collects unpublished results, solver-design sketches, and
research directions that extend the thesis, adapted from the research
proposals in my second-year report. We begin with the directions we consider
most fertile---each backed by a concrete unpublished result or a concrete
algorithmic plan---and close with briefer sketches and a broader outlook.

\section{Nonlinear Bitvector Arithmetic via the 2-adic Integers}
\label{sec:fw:2adic}

The solvers of \cref{ch:single-width-bv,ch:multi-width-bv} decide the
\emph{linear-bitwise} fragment of the parametric bitvector theory. We have an
unpublished result showing that \emph{nonlinear} arithmetic in the parametric
bitvector theory is decidable, by a reduction to the first-order theory of
the 2-adic integers. On the other hand, our unpublished results also show
that integrating \emph{any} bitwise operation into such a theory immediately
leads to undecidability.\footnote{%
\url{https://pixel-druid.com/articles/non-linear-theory-of-2-adics-does-not-mix-with-bitwise-operations}}
This implies that the nonlinear arithmetic fragment of the parametric
bitvector theory is (maximally) decidable: to write incomplete solvers for
the full theory, one must first build a complete solver for the nonlinear
arithmetic fragment, which can then be mixed with bitwise operations via
theory combination.

Toward this, we propose implementing a fast cylindrical algebraic
decomposition (CAD) solver for the 2-adics. While the decidability of the
full first-order theory of the 2-adics has been known in theory since the
1960s~\citep{ax1966diophantine}, that result was not algorithmic. A CAD-style
algorithm based on cell decomposition does
exist~\citep{quantifier-elimination-padic}, but, to the best of our
knowledge, no implementation of it does. We therefore propose to implement
such a solver in C++ or Rust, leveraging existing libraries such as
FLINT~\citep{hart2013flint} for polynomial manipulation and root finding.
This also sets the stage for a verified reimplementation in Lean, extending
the reach of the verified solvers of this thesis beyond the linear-bitwise
fragment.
% See also: Solving Satisfiability of Polynomial Formulas By Sample-Cell
% Projection, https://arxiv.org/pdf/2003.00409

\section{Ackermannization, Incremental Solving, and the Theory of Arrays}
\label{sec:fw:ackermannization}

Adding \emph{incremental solving} to \bvdecide{} would enable much better
variants of ackermannization, thereby unlocking support for the theory of
arrays. Together with the floating-point support of \cref{ch:floating-point},
this would bring the verified solver to theory parity with
Bitwuzla~\citep{niemetz2023bitwuzla}
($\qfbv{} + \texttt{QF\_FP} + \texttt{QF\_UF}$), which suffices to model
almost all behaviour of realistic production languages such as C and C++.
Industry formalization efforts would be a natural testing ground for
evaluating whether such an extended solver can handle real verification
workloads. Bitwuzla's ackermannization and multiplication-abstraction
machinery is a natural starting point for a mechanization.

\section{Average-Case Floating-Point Error Analysis via \texorpdfstring{$\delta$}{delta}-Satisfiability}
\label{sec:fw:deltasat}

Worst-case floating-point error bounds are a key ingredient for optimizing
machine-learning kernels, yet they are often too pessimistic for
machine-learning workloads. We propose to compute tighter,
\emph{average-case} error bounds by building on the dReal
solver~\citep{gao2013dreal}, which decides \emph{bounded} quantified problems
over the reals. The key idea of dReal is that for Lipschitz-continuous
functions over bounded domains, one can compute a (uniform) discretization of
the domain that is sufficiently fine to decide satisfiability up to a given
$\delta$; this is $\delta$-satisfiability, solved by a CDCL-style algorithm.
The key insight for our proposal is that average-case pointwise error
analysis for machine-learning kernels reduces to bounded quantified formulas
over the reals extended with integration. A recent solver,
$\int$dreal~\citep{rivera2025checking}, supports exactly this theory, and we
plan to adapt it to floating-point error analysis.

\section{Verified Solving for Tensor Reshaping and Reindexing}
\label{sec:fw:tensor}

Tensor equivalence under reshaping and reindexing is a key problem in
optimizing machine-learning kernels. We believe a sound and complete solver
for this problem is within reach, by reducing it to a combination of
(a)~the word problem in the symmetric group, for reindexing, and
(b)~parametric bitvector solving, for reshaping.

It may appear curious to use a bitvector solver for reshaping. However,
recent work on tensor code generation~\citep{zhou2025linear} shows that most
tensor layout problems use powers of two, and thus reduce to linear algebra
over \Ftwo{}---which is efficiently encoded as parametric bitvector
problems of exactly the shape our solvers from \cref{ch:multi-width-bv}
handle. While both component theories have known solvers, to the best of our
knowledge there is no solver for their \emph{combination}, and no verified
solver for the word problem in the symmetric group. A beautiful theory of
strong generating sets permits fast decision procedures not just for the word
problem but for almost all structural properties of subgroups of the
symmetric group~\citep{seress2003permutation}. We propose to build a verified
solver based on strong generating sets and combine it with our parametric
bitvector solver. Insights from this theory combination could further lead to
solvers for \emph{parametric bitvectors with permutations}---a powerful
theory for reasoning about vectorized code that uses the complex shuffle
operations of instruction sets such as AVX-512.

\section{Better Model Guided Generalization for Exists Forall in QF\_BV}
\label{sec:fw:modelguided}

Hydra algorithm is naive, find better model guided generalization,
using ideas from virtual substitution for the reals.

\subsection{Virtual Substitution}
\subsubsection{Thirty Years of Virtual Substitution}
\subsubsection{New Concepts for Real Quantifier Elimination by Virtual Substitution}
\subsubsection{Simplification of Quantifier-free Formulae over Ordered Fields}
% \subsubsection{A continuous, constructive solution to Hilbert's 17th problem}

\section{Reducing Bounded NLIA to QF\_BV}
\label{sec:fw:nlia}

\section{Extending MBA to Inequational Fragment plus Disjunction}
\label{sec:fw:mba}

- Unsigned Inequations can be handled using the same idea, see hacker's delight, Ch2.2
- Can disjunction be handled the same way? unclear

\section{Multiple-Widths in Width Generic Automataion}
\label{sec:fw:multiwidth}

Key idea: If we have widths $x : w$ and $ y : w+k$, note that this really means that
if we are at width $w$, then the high bits of $x$ are zero when treated as a transducer.
So this should allow us to exploit the structure to decide properties.

\section{Exponentition, Modulus, and Monus}
\label{sec:fw:exponentiation}

See that if we have $2^{w-k}$ and $2^{w - l}$, we really can work in a regime where $w \geq k, w \geq l$,
and just bitblast the case where $w \leq max(k, l)$.
WLOG, $k \geq l$, then we can set $k = l + \delta$, and pick $w' \geq l$.

\section{Broader Outlook: AI for Formal Verification}
\label{sec:fw:ai}

Beyond decision procedures themselves, two directions apply the
infrastructure of this thesis to machine-assisted formalization.

\paragraph{Large-scale auto-formalization of software engineering theories.}
Formal libraries for bitvector and floating-point theories remain sparse,
owing to a lack of both training data (such as textbooks) and domain-specific
tactics that work well over these rich domains. The tactic infrastructure and
domain knowledge developed in this thesis provide the foundation needed to
bootstrap large-scale formalization efforts for these theories, given
sufficient models and compute.

\paragraph{Teaching models to metaprogram in Lean.}
A principled tactic in Lean comprises (a)~a decision procedure for the
domain, and (b)~the metaprogramming glue connecting it to the user frontend
(proof by reflection, \texttt{simp}, \texttt{grind}, simprocs, and so on).
The metaprogramming API is complex, underdocumented, and contains subtle
invariants that are easy to break; if AI models are to create their own
tactics, they must master this metaprogramming layer. We would like to
explore whether cleaner abstractions---analogous to how \texttt{SymM} inside
\texttt{GrindM} abstracts away the metaprogramming machinery for symbolic
simulation---can be designed for proof by reflection, making the API
tractable for both humans and models.
